August 3, 2026

Editor-in-Chief
IEEE Sensors Journal
IEEE Sensors Council

SUBJECT: Submission of Manuscript P6 — Non-Semantic Acoustic Sentinel for Classroom Noise Anomaly Detection

Dear Editor-in-Chief and Associate Editors,

We are pleased to submit our original research manuscript titled "Non-Semantic Acoustic Sentinel for Classroom Noise Anomaly Detection" for publication consideration in the IEEE Sensors Journal.

1. SUBSYSTEM BACKGROUND AND MOTIVATION
Acoustic sensing in educational environments is vital for detecting physical distress or safety disturbances. However, continuous microphone recording violates participant privacy and violates eavesdropping regulations. Transcribing speech or storing raw voice waveforms creates high biometric exposure. This paper introduces a non-semantic acoustic sensing pipeline that classifies ambient safety events without capturing intelligible speech.

2. CORE TECHNICAL CONTRIBUTIONS AND NOVELTY
  - Non-Semantic Spectral Mapping: Processes 100ms PCM audio buffers using Fast Fourier Transform (FFT) equations to extract Spectral Centroid, Zero-Crossing Rate (ZCR), and Spectral Flux.
  - Volatile Memory Deallocation: Immediately overwrites raw PCM audio buffers in RAM, proving that linguistic speech reconstruction from the stored frequency feature vectors is impossible.
  - Context-Aware Acoustic Thresholding: Calibrates noise sensitivity thresholds dynamically based on timetable modes (Exam vs. Break), preventing false alarms during normal social interactions.

3. SUITABILITY FOR IEEE SENSORS JOURNAL
This manuscript fits the IEEE Sensors Journal's focus on acoustic sensing applications, real-time signal processing, privacy-preserving sensor nodes, and environmental anomaly detection.

4. ETHICS AND ORIGINALITY DECLARATIONS
This work represents original research, is not under consideration elsewhere, and strictly adheres to institutional acoustic privacy standards.

Sincerely,

Polisetti Narendra (Corresponding Author)
Department of Computer Science & Engineering
Swarnandhra College of Engineering and Technology (Autonomous)
Seetharampuram, Narasapur-534280, Andhra Pradesh, India
Email: narendraa1996@gmail.com
